\documentclass[12pt,a4paper]{article}
\usepackage[utf8]{inputenc}
\usepackage[T1]{fontenc}
\usepackage{geometry}
\geometry{margin=1in}
\usepackage{hyperref}
\usepackage{listings}
\usepackage{xcolor}
\usepackage{graphicx}
\usepackage{titlesec}
\usepackage{enumitem}

\definecolor{codegreen}{rgb}{0,0.6,0}
\definecolor{codegray}{rgb}{0.5,0.5,0.5}
\definecolor{codepurple}{rgb}{0.58,0,0.82}
\definecolor{backcolour}{rgb}{0.95,0.95,0.92}

\lstdefinestyle{mystyle}{
    backgroundcolor=\color{backcolour},
    commentstyle=\color{codegreen},
    keywordstyle=\color{magenta},
    numberstyle=\tiny\color{codegray},
    stringstyle=\color{codepurple},
    basicstyle=\ttfamily\footnotesize,
    breakatwhitespace=false,
    breaklines=true,
    captionpos=b,
    keepspaces=true,
    numbers=left,
    numbersep=5pt,
    showspaces=false,
    showstringspaces=false,
    showtabs=false,
    tabsize=2
}

\lstset{style=mystyle}

\title{
    \vspace*{2cm}
    \textbf{Python GUI Calculator} \\
    \large Modern Implementation using Clean Architecture and TDD
}
\author{
    \textbf{Tal Fiskus} \\
    BIU Reinforcement Learning (RL) Course no. 896873 \\
    Bar Ilan University
}
\date{March 18, 2026}

\begin{document}

\maketitle
\newpage

\section*{Abstract}
\addcontentsline{toc}{section}{Abstract}
This report details the design and implementation of a desktop-based Graphical User Interface (GUI) calculator developed using Python. The primary objective was to apply professional software engineering standards, specifically Clean Architecture and Test-Driven Development (TDD). The resulting application features a high-precision mathematical engine utilizing the \texttt{decimal} module, a configuration-driven UI layer, and a robust software development kit (SDK) to bridge the two. Validation via automated testing achieved over 90\% code coverage for the core logic, ensuring reliability, accuracy, and maintainability.

\section{Introduction}
\subsection{Problem Statement}
Basic calculator applications often suffer from a lack of modularity, hard-coded UI elements, and precision issues inherent in standard binary floating-point arithmetic. Most elementary implementations tightly couple the logic with the display, making them difficult to maintain, test, or extend for scientific purposes.

\subsection{Motivation}
The motivation for this project was to demonstrate that even a seemingly simple application can benefit from rigorous architectural patterns. By employing \textit{Clean Architecture}, the project aims to isolate the mathematical "domain" from implementation details like the UI framework (Tkinter). This ensures that the core logic can be tested independently of the interface and that the UI can be replaced or updated (e.g., switching from Tkinter to PyQt) without affecting the underlying engine.

\section{Methodology}
\subsection{Architectural Decisions: Clean Architecture}
The project is structured into three distinct layers to ensure separation of concerns:
\begin{itemize}
    \item \textbf{Domain Layer (Core):} Contains pure Python business logic (the mathematical engine). It has zero dependencies on external libraries or other layers.
    \item \textbf{Application Layer (SDK):} Orchestrates the flow of data between the UI and the Core. It handles input sanitization, error wrapping, and output formatting.
    \item \textbf{Infrastructure/UI Layer:} Implements the graphical interface using the Tkinter library. It reads all stylistic and structural data from an external configuration file.
\end{itemize}

\subsection{Design Patterns}
\begin{itemize}
    \item \textbf{SDK/Facade Pattern:} The \texttt{CalculatorSDK} provides a simplified interface for the UI, hiding the complexity of the internal state machine.
    \item \textbf{Configuration Pattern:} All "magic numbers," strings, and color hex codes are moved to a JSON configuration file to achieve zero hard-coding.
    \item \textbf{TDD (Test-Driven Development):} All engine features were implemented only after a corresponding failing test was written, ensuring 100\% behavioral alignment with the PRD.
\end{itemize}

\section{Implementation}
\subsection{Code Structure}
The project follows a modular directory structure:
\begin{itemize}
    \item \texttt{src/core/engine.py}: The mathematical "brain" managing operands and operators.
    \item \texttt{src/sdk/calculator\_sdk.py}: The bridge layer handling display limits and formatting.
    \item \texttt{src/ui/}: Divided into \texttt{components.py} (widgets) and \texttt{main\_window.py} (layout).
    \item \texttt{config/settings.json}: The single source of truth for UI configuration.
\end{itemize}

\subsection{Key Algorithms}
The \texttt{CalculatorEngine} utilizes Python's \texttt{Decimal} module to handle arithmetic. Unlike standard floats, \texttt{Decimal} provides exact decimal representation, crucial for avoiding errors like $0.1 + 0.2 \approx 0.30000000000000004$. The engine implements a state-based logic to handle sequential operations (e.g., $5 + 5 + 5$) correctly.

\section{Results \& Testing}
\subsection{Test Coverage}
Implementation was validated using the \texttt{pytest} framework. 18 comprehensive test cases were written, covering:
\begin{itemize}
    \item Basic arithmetic (+, -, *, /).
    \item Sequential operations and equals-key behavior.
    \item Edge cases: Division by zero, multiple decimal points, and large number overflow.
\end{itemize}

\subsection{Performance Metrics}
\begin{itemize}
    \item \textbf{Core Engine Coverage:} 92\%
    \item \textbf{Overall Project Coverage:} 87\%
    \item \textbf{Linter Compliance:} 100\% PEP 8 compliance via \texttt{flake8}.
    \item \textbf{UI Responsiveness:} Sub-millisecond response time for all button clicks.
\end{itemize}

\section{Conclusion}
\subsection{Achievements}
The project successfully delivered a functional, aesthetically pleasing, and technically sound calculator application. It proves that Clean Architecture is highly effective for isolating logic, and TDD significantly reduces the time spent on manual debugging.

\subsection{Future Improvements}
\begin{itemize}
    \item \textbf{Keyboard Integration:} Support for physical numpad input mapping.
    \item \textbf{Scientific Mode:} Extending the engine with trigonometric and logarithmic functions.
    \item \textbf{History Tracking:} Implementing a side panel to review past calculation sequences.
\end{itemize}

\newpage
\appendix
\section{Appendices}

\subsection{Configuration Snippet}
\begin{lstlisting}[language=JSON, caption=config/settings.json]
{
    "window": { "title": "Python GUI Calculator", "width": 400, "height": 500 },
    "display": { "max_chars": 15, "bg": "#f0f0f0" },
    "buttons": {
        "operators": { "bg": "#ff9500", "fg": "#ffffff", "hover": "#e08400" }
    }
}
\end{lstlisting}

\subsection{Core Math Logic}
\begin{lstlisting}[language=Python, caption=src/core/engine.py (Simplified)]
from decimal import Decimal

def calculate(self):
    try:
        val1 = Decimal(str(self.previous_input))
        val2 = Decimal(str(self.current_input))
        if self.active_operator == "+":
            result = val1 + val2
        # ... logic for other operators
        self.current_input = str(result.normalize())
    except ZeroDivisionError:
        self.current_input = "Error"
\end{lstlisting}

\end{document}
